\section{Experimental Design}
\label{sec:protocol}

\paragraph{Data and split.}
We use NIH ChestX-ray14 \citep{wang2017chestxray}, whose findings are mined from radiology reports. A deterministic hash of patient identity assigns rows to disjoint train, validation, and test partitions. The registered manifest contains 18,212 training images and 5,343 test images from 1,659 test patients. The three cells pair LLaVA-1.5-7B with Effusion and Edema, and Qwen2.5-VL-7B with Effusion. All reported uncertainty clusters by patient. Prompts and exact run identifiers appear in Appendix~\ref{app:protocol}.

\paragraph{Models and consumed loci.}
For LLaVA-1.5-7B \citep{liu2024improved} and Qwen2.5-VL-7B \citep{bai2025qwen25vl}, we use the final visual block consumed by the connector or merger. Synthetic hook tests establish downstream-path membership, a bitwise $\alpha=0$ no-op, and changed downstream representations and logits at nonzero dose. Probe and intervention therefore share one consumed representation; exact modules appear in Appendix~\ref{app:protocol}.

\paragraph{Controlled decodability.}
For image $i$, mean pooling over the final-block visual tokens yields $\hvec_i\in\R^D$. A fixed seed-0 Gaussian matrix $R\in\R^{D\times512}$ projects every model to a common readout dimension; featurewise scale $\boldsymbol{s}$ is fitted on training rows only. We then fit logistic regression with $C=1$, maximum 2,000 iterations, and seed 0. The real-label AUROC is compared with 20 control AUROCs. Each control seed maps the same recurring input type---AP view indicator, sex indicator, and age decade, giving 39 strata---to a random label. Selectivity is
\begin{equation}
S = \operatorname{AUROC}(y,\hat y)-\frac{1}{20}\sum_{s=0}^{19}\operatorname{AUROC}(\tilde y_s,\hat y_s).
\label{eq:selectivity}
\end{equation}
We report the real AUROC, control distribution, and $S$ with 2,000 patient-cluster bootstrap resamples. A positive interval for $S$ means that the target is more readable than same-capacity nuisance-type label maps. It does not yet imply downstream use. Paired contrasts reuse identical ordered test rows and bootstrap draws.

\paragraph{Probe-normal intervention.}
Let $\boldsymbol\beta_c$ be the projected-space logistic coefficient for the positive class of concept $c$. We lift it to the consumed $D$-dimensional representation and normalize it exactly as
\begin{equation}
\hat\wvec_c = \frac{R\left(\boldsymbol\beta_c\oslash\max(\boldsymbol{s},10^{-8})\right)}{\left\lVert R\left(\boldsymbol\beta_c\oslash\max(\boldsymbol{s},10^{-8})\right)\right\rVert_2}.
\label{eq:lift}
\end{equation}
The positive-class coefficient fixes the direction's clinical pole. At every visual token $t$, the registered intervention is
\begin{equation}
\hvec'_{it}=\hvec_{it}+\alpha\lVert\hvec_{it}\rVert_2\hat\wvec_c.
\label{eq:intervention}
\end{equation}
Here the sign of $\alpha$ selects the positive or negative clinical pole, while $|\alpha|$ is the perturbation norm as a fraction of each token's activation norm. The endpoint is the paired change in mean $P(yes)$ from the common unsteered baseline. Original cells use concept doses $\{-1,-0.5,-0.25,-0.1,0,0.1,0.25,0.5,1\}$ on 200 held-out test images. Controls are evaluated at the six nonzero doses excluding $\pm0.1$: 20 isotropic random unit directions, one sham formed by a fixed coordinate permutation of $\hat\wvec_c$ and renormalization, and fixed alternative clinical probe normals. Every control receives the same signed dose on the same rows. Dose monotonicity over $[-0.5,0.5]$ is descriptive.

\paragraph{Original-cell decision rule.}
For each controlled nonzero dose, the concept-consistent change is $\operatorname{sign}(\alpha)\Delta(\alpha)$: positive steering should increase $P(yes)$, and negative steering should decrease it. The primary effect is the maximum concept-consistent change over those doses. A cell clears the generic intervention gate only if this effect leaves the same-dose random 5th--95th percentile interval in the concept-consistent direction and exceeds the maximum absolute sham effect. The random p95 is a registered reference gate over 20 directions, not a conventional 5\% hypothesis test. The conjunction is fixed across the three cells.

\paragraph{Direction-specificity supplement.}
The initial experiment prespecified the random and sham comparisons; its clinical-direction analysis was descriptive. The Nodule response motivated a separately registered supplement at the locked Qwen dose, $\alpha=+0.25$. We select one row by label-blind hash for each of 400 ChestX-ray14 test patients absent from the original 200-image intervention cohort. The semantic family contains all five non-target normals from the pre-existing Qwen direction bundle---Atelectasis, Pneumothorax, Cardiomegaly, Mass, and Nodule---rather than directions chosen from the supplement outcome. Each uses the same training rows, projection, scaling, classifier, lifting rule, and normalization as Effusion. Together with a common baseline, 20 random directions, and the sham, the grid produces 28 configurations and 11,200 per-image outcomes; exact row selection appears in Appendix~\ref{app:protocol}.

For direction $d$, let $\Delta_d$ be its paired mean change from baseline. The familywise statistic is
\begin{equation}
M=\DeltaEff-\max_{d\in\mathcal D_{\mathrm{clinical}}}\Delta_d.
\label{eq:margin}
\end{equation}
The maximum is recomputed inside each of 5,000 patient-bootstrap draws. Thus $M>0$ means Effusion beats every member of the fixed clinical family, while $M<0$ means at least one alternative is stronger. Direction specificity requires a positive Effusion effect above random p95 and absolute sham, together with a one-sided 95\% lower bound for $M$ above zero.
